\documentclass[10pt,a4paper]{amsart}
\usepackage[margin=19mm]{geometry}
\usepackage{amsmath,amssymb,mathtools}
\usepackage[scr=rsfs,cal=euler]{mathalfa}
\usepackage{xfrac}
\usepackage[unicode,hidelinks,bookmarks=false]{hyperref}
\newcommand{\ie}{i.e.\ }
\newcommand{\cf}{cf.\ }
\newcommand{\resp}{respectively\ }
\newcommand{\Subsection}{\subsection}
\providecommand{\nobreakdash}{\nobreak\hbox{-}}
\setlength{\emergencystretch}{2em}
\multlinegap0pt
\usepackage{amssymb}
\newtheorem{corollary}{Corollary}
\newtheorem{proposition}{Proposition}
\newtheorem{theorem}{Theorem}
\title{\textsf{A Constructive Version of Ekeland's Variational Principle}}
\author{Douglas S. Bridges}
\date{Source: arXiv:2606.10339v2 (CC BY 4.0)}
\begin{document}
\maketitle

\noindent\emph{Source and licence.} \textsf{A Constructive Version of Ekeland's Variational Principle}. Version arXiv:2606.10339v2 (CC BY 4.0), distributed by arXiv under the Creative Commons Attribution 4.0 International licence, http://creativecommons.org/licenses/by/4.0/, as recorded in the arXiv licence metadata for this exact version and shown on the abstract page https://arxiv.org/abs/2606.10339v2. Full licence notice: \texttt{licenses/arxiv-2606.10339-CC-BY-4.0.txt}.\par\smallskip

\noindent\textit{Editorial scope.} Counted unit: the article's theorem Ekeland (source0.tex lines 505--517). The article's complete original proof of it is delivered with it (source0.tex lines 519--650, one \texttt{proof} environment of 132 lines). No mathematics is written, repaired or re-derived: the statement and the proof are the article's own lines, in the article's own notation, and no retained line is substituted or re-flowed.  Retained uncounted as the source-local context the proof needs: lines 375--381; lines 419--473.  Nothing was omitted from any retained span, so the package carries no omission record.  No retained line contains a citation, so the wrapper carries no bibliography and no bibliography entry.  No retained line contains a figure environment, an image-inclusion command or drawing code, so no image is delivered and the package declares no asset dependency.  Editorial formatting only, in addition: an AMS \texttt{amsart} preamble that restates the article's own declarations and macros used by the retained lines, namely the environments \texttt{corollary}, \texttt{proposition}, \texttt{theorem} on separate counters, together with the standard packages those lines need.  The retained environments print in this document as Corollary 1, Proposition 1, Theorem 1 in order of appearance, whereas the article prints the same items with the numbering of its own full document.  The complete native source of the article as distributed in the arXiv e-print for this exact version is bundled unchanged as \texttt{sources/202465/source0.tex}.
\begin{corollary}
\label{may25c3}Let $f$ be a continuous mapping of a locally compact metric
space $X$ into $\mathbb{R}$ such that $\inf f$ exists. Let $a\in X$ and $c>0$.
Then the infimum of the mapping $g:x\rightsquigarrow f(x)+c\rho(x,a)$ on $X$
exists, and $\mathsf{L}(g,\lambda)$ is compact for all but countably many
$\lambda>\inf g$.
\end{corollary}

\begin{proposition}
\label{may04p1}Let $X$ be a metric space, and $\left(  K_{n}\right)
_{n\geqslant1}$ a sequence of compact subsets of $X$ that

\begin{itemize}
\item[\emph{(a)}] is decreasing---that is, $K_{n+1}\subset K_{n}$ for each
$n$--- and\footnote{%
%TCIMACRO{\TeXButton{sf}{\normalfont\sf}}%
%BeginExpansion
\normalfont\sf
%EndExpansion
Dropping hypothesis (b) from Proposition \ref{may04p1}, and using the
classical---but not constructive---sequential compactness of the sets $K_{n}$,
we can easily prove that, simply under the hypothesis (a), $%
%TCIMACRO{\tbigcap _{n\geq1}}%
%BeginExpansion
{\textstyle\bigcap_{n\geq1}}
%EndExpansion
K_{n}$ is inhabited. Here, however, is a Brouwerian example showing that we
cannot prove this constructively. Let $\left(  a_{n}\right)  _{n\geq1}$ be a
binary sequence with at most one term equal to $1$ (more formally, $a_{m}%
a_{n}=0$ whenever $m\neq n$). If $a_{k}=0$ for all $k\leq n$, set
$K_{n}=\left\{  0,1\right\}  $; if $a_{n}=1$, then set $K_{n}=\left\{
0\right\}  $ if $n$ is even, and $K_{n}=\{1\}$ if $n$ is odd. Then $\left(
K_{n}\right)  _{n\geq1}$ is a decreasing sequence of compact subsets of
$\left\{  0,1\right\}  $. If there exists $\xi$ in $%
%TCIMACRO{\tbigcap _{n\geq1}}%
%BeginExpansion
{\textstyle\bigcap_{n\geq1}}
%EndExpansion
K_{n}$, then either $\xi>0$ or $\xi<1$. In the first case, $a_{n}=0$ for all
even $n$, and in the second, $a_{n}=0$ for all odd $n$. Thus the
aforementioned classical result implies the essentially nonconstructive
principle \textsf{LLPO: \ }If $\left(  a_{n}\right)  _{n\geq1}$ is a binary
sequence with at most one term equal to $1$, then either all even-indexed
terms $a_{n}$ equal $0$, or else all odd-indexed terms equal $0$.}

\item[\emph{(b)}] such that $\mathsf{diam}(K_{n})\rightarrow0$ as
$n\rightarrow\infty$.
\end{itemize}

%

%TCIMACRO{\TeXButton{noindent}{\noindent}}%
%BeginExpansion
\noindent
%EndExpansion
For each $n$ let $x_{n}\in K_{n}$. Then $\left(  x_{n}\right)  _{n\geqslant1}$
converges to the unique point $\xi$ in $%
%TCIMACRO{\tbigcap _{n\geqslant1}}%
%BeginExpansion
{\textstyle\bigcap_{n\geqslant1}}
%EndExpansion
K_{n}$.
\end{proposition}

\begin{theorem}
\label{Ekeland}Let $f$ be a continuous real-valued mapping on a locally
compact metric space $X$, such that $\inf f$ exists. Let $x_{0}\in X$, $c>0,$
and $\varepsilon>0$. Then there exists $v\in X$ such that%
\begin{equation}
f(v)+c\rho\left(  v,x_{0}\right)  \leqslant f\left(  x_{0}\right)  \label{1}%
\end{equation}
and%
\begin{equation}
f\left(  v\right)  <f(x)+c\rho\left(  v,x\right)  +\varepsilon. \label{2}%
\end{equation}
for all $x\neq v$ in$\ X$.
\end{theorem}

\begin{proof}
Define the mapping $g_{0}:X\rightarrow\mathbb{R}$ by $g_{0}(x)\equiv
f(x)+c\rho(x,x_{0})$. Then $\inf g_{0}$ exists, by Corollary \ref{may25c3}. We
may assume that $\inf f=0$ and therefore $\inf g_{0}$ is positive. By
Corollary \ref{may25c3}, there exists $\varepsilon_{0}\ $with $0<\varepsilon
_{0}<\frac{1}{2}\min\left\{  \varepsilon,f(x_{0})\right\}  $ and
$f(x_{0})-\varepsilon_{0}\neq\inf g_{0}$, such that if $f(x_{0})-\varepsilon
_{0}>\inf g_{0}$, then $\mathsf{L}(g_{0},f(x_{0})-\varepsilon_{0})$ is
compact. If $f(x_{0})-\varepsilon_{0}<\inf g_{0}$, then $\mathsf{L}%
(g_{0},f(x_{0})-\varepsilon_{0})$ is empty, in which case for all $x\in X$,
\[
f(x_{0})\leqslant f(x)+c\rho(x,x_{0})+\varepsilon_{0}<f(x)+c\rho
(x,x_{0})+\varepsilon
\]
and $x_{0}$ itself satisfies the desired conditions for $v$. Thus we may
assume that $\mathsf{L}(g_{0},f(x_{0})-\varepsilon_{0})$\ is compact. Setting
$\lambda_{0}=0$ and $K_{0}=\mathsf{L}(g_{0},f(x_{0})-\varepsilon_{0})$, we
construct an increasing binary sequence $\left(  \lambda_{n}\right)
_{n\geqslant0}$; a sequence $\left(  g_{n}\right)  _{n\geqslant0} $ of
continuous real-valued functions on $X$; a sequence $\left(  K_{n}\right)
_{n\geqslant0}$ of compact subsets of $X$; a sequence $\left(  \varepsilon
_{n}\right)  _{n\geqslant0}$ of positive numbers; and a sequence $\left(
x_{n}\right)  _{n\geqslant0}$ of points of $X$, such that the following hold
for each applicable $n$.

\begin{itemize}
\item[(a)] $0<\varepsilon_{n}<\frac{1}{2}\min\left\{  \varepsilon
,f(x_{n})\right\}  $.

\item[(b)] If $\lambda_{n}=0$, then $g_{n}(x)=f(x)+c\rho(x,x_{n})$,
$f(x_{n})-\varepsilon_{n}>\inf g_{n}$,$\ K_{n}\equiv\mathsf{L}(g_{n}%
,f(x_{n})-\varepsilon_{n})$ is compact, $x_{n+1}\in K_{n}\,$, and
$f(x_{n+1})<\inf f(K_{n})+2^{-n-1}$.

\item[(c)] If $\lambda_{n}=1-\lambda_{n-1}$, then $g_{n}(x)=f(x)+c\rho
(x,x_{n})$, $f(x_{n})-\varepsilon_{n}<\inf g_{n}$, and for all $m\geqslant n $
we have $\lambda_{m}=1$, $g_{m}=g_{n}$, $K_{m}=\left\{  x_{n}\right\}  $, and
$x_{m}=x_{n}$.

\item[(d)] If $n\geqslant1$, then $K_{n}\subset K_{n-1}$.
\end{itemize}

%

%TCIMACRO{\TeXButton{noindent}{\noindent}}%
%BeginExpansion
\noindent
%EndExpansion
Suppose that for some $n\geqslant0$ we have constructed $\lambda_{n-1}%
,g_{n-1},K_{n-1},\varepsilon_{n-1},$ and $x_{n-1}$ with the applicable
properties. If $\lambda_{n-1}=0$, choose $x_{n}\in K_{n-1}$ such that
$f(x_{n})<\inf f(K_{n-1})+2^{-n}$. By Corollary \ref{may25c3}, $\inf g_{n}$
exists, as does $\varepsilon_{n}$ with $0<\varepsilon_{n}<\frac{1}{2}%
\min\left\{  \varepsilon,f(x_{n})\right\}  $ and $f(x_{n})-\varepsilon_{n}%
\neq\inf g_{n}$, such that if $f(x_{n})-\varepsilon_{n}>\inf g_{n}$, then
$\mathsf{L}(g_{n},f(x_{n})-\varepsilon_{n})\ $is compact. Either
$f(x_{n})-\varepsilon_{n}<\inf g_{n}$, in which case $\mathsf{L}(g_{n}%
,f(x_{n})-\varepsilon_{n})$ is empty and we set $\lambda_{n}=1$ and
$K_{n}=\left\{  x_{n}\right\}  $; or else $f(x_{n})-\varepsilon_{n}>\inf
g_{n}$, in which case we set $\lambda_{n}=0$ and $K_{n}=\mathsf{L}%
(g_{n},f(x_{n})-\varepsilon_{n})$. This disposes of the alternative
$\lambda_{n-1}=0$. If, on the other hand, $\lambda_{n-1}=1$, we set
$\lambda_{n}=1$, $g_{n}=g_{n-1}$, $K_{n}=\left\{  x_{n-1}\right\}  $,
$\varepsilon_{n}=\varepsilon_{n-1}$, and $x_{n}=x_{n-1}$. This completes the
inductive construction itself. Clearly, (a) and (b) are satisfied. It is
straightforward to prove (c), so it remains to prove (d). To that end, if
$\lambda_{n}=0$, then for each $x\in K_{n}$ we have
\begin{align*}
f\left(  x\right)  +c\rho\left(  x,x_{n-1}\right)   &  \leqslant f\left(
x\right)  +c\rho\left(  x,x_{n}\right)  +c\rho\left(  x_{n},x_{n-1}\right) \\
&  \leqslant f\left(  x_{n}\right)  -\varepsilon_{n}+c\rho\left(
x_{n},x_{n-1}\right) \\
&  \leqslant f\left(  x_{n-1}\right)  -\varepsilon_{n}-\varepsilon
_{n-1}\ \ \text{(as }x_{n}\in K_{n-1}\text{)}\\
&  <f\left(  x_{n-1}\right)  -\varepsilon_{n-1}%
\end{align*}
and therefore $x\in K_{n-1}$; whence $K_{n}\subset K_{n-1}$. If $\lambda
_{n}=1-\lambda_{n-1}$, then $K_{n}=\left\{  x_{n}\right\}  \subset K_{n-1}$,
by (c)~and (b); and if $\lambda_{n-1}=1$, then $K_{n}=K_{n-1}$, by (c). This
completes the proof of (d).

If $n\geqslant1$ and $\lambda_{n}=0$, then for each $x\in K_{n}$,%
\[
\inf f(K_{n})+c\rho\left(  x,x_{n}\right)  \leqslant f\left(  x\right)
+c\rho\left(  x,x_{n}\right)  \leqslant f\left(  x_{n}\right)  -\varepsilon
_{n},
\]
so, noting that $\inf f(K_{n})\geqslant\inf f(K_{n-1})$, we have
\[
\rho\left(  x,x_{n}\right)  \leqslant c^{-1}(f\left(  x_{n}\right)  -\inf
f(K_{n})-\varepsilon_{n})<c^{-1}(f\left(  x_{n}\right)  -\inf f(K_{n-1}%
))<\frac{1}{2^{n}c}.
\]
It follows from this that $\mathsf{diam}(K_{n})<1/2^{n-1}c$ for all $n$. Hence
$\left(  K_{n}\right)  _{n\geqslant0}$ is a descending sequence of compact
sets whose diameters tend to $0$. By Proposition \ref{may04p1}, $\left(
x_{n}\right)  _{n\geqslant1}$ converges to the unique point$\ v~$of $%
%TCIMACRO{\tbigcap _{n\geqslant0}}%
%BeginExpansion
{\textstyle\bigcap_{n\geqslant0}}
%EndExpansion
K_{n}$. Since $v\in K_{0}$, statement (\ref{1}) clearly holds. On the other
hand, given $x\in X$ with $x\neq v$, we have either (\ref{2}) or%
\begin{equation}
f\left(  x\right)  +c\rho\left(  x,v\right)  <f\left(  v\right)
-\frac{\varepsilon}{2}. \label{06a}%
\end{equation}
To complete the proof, it will suffice to rule out the latter alternative.
Assuming (\ref{06a}), if $\lambda_{n}=1-\lambda_{n-1}$ we have $\mathsf{L}%
(g_{n},f(x_{n})-\varepsilon_{n})=\varnothing$ and $x_{m}=x_{n}$ for all
$m\geqslant n$. Hence $v=x_{n}$ and therefore%
\[
f\left(  x\right)  +c\rho\left(  x,v\right)  =f\left(  x\right)  +c\rho\left(
x,x_{n}\right)  \geqslant f\left(  x_{n}\right)  -\varepsilon_{n}\geqslant
f\left(  v\right)  -\frac{\varepsilon}{2},
\]
contradicting (\ref{06a}). Thus, recalling that $\lambda_{0}=0$, for all $n$
we must have $\lambda_{n}=0$ and $K_{n}$ equal to the compact set$\ \mathsf{L}%
(g_{n},f(x_{n})-\varepsilon_{n})$. By (\ref{06a}) and the continuity of $f $
and $\rho$, for all sufficiently large $n$ we have%
\[
f(x)+c\rho(x,x_{n})+\varepsilon_{n}<f(x)+c\rho(x,x_{n})+\frac{\varepsilon}%
{2}<f(x_{n});
\]
whence $x\in K_{n}$ and therefore%
\[
\rho(x_{n},x)\leqslant\mathsf{diam}(K_{n})<\frac{1}{2^{n-1}c}\rightarrow
0\text{ as }n\rightarrow\infty\,.
\]
Thus $x=\lim_{n\rightarrow\infty}x_{n+1}=v$, a contradiction which ensures
that the alternative (\ref{06a}) is ruled out.
\end{proof}

\end{document}
